\documentclass[11pt]{article}

\usepackage{amsmath,amssymb}
\usepackage{xurl}
\usepackage[hidelinks]{hyperref}
\setlength{\emergencystretch}{2em}

\input{command}

\title{Dimension Divisibility in the Classical Test Instantiation of the
	Pauli Basis Analysis}
\author{}
\date{2026-08-30}

\begin{document}
\maketitle

\section{At a glance}

\begin{itemize}
  \item \textbf{Difficulty: medium.}  Two related matters are recorded.
  First, the soundness analysis of the Pauli basis test invokes the
  classical low individual degree game at dimension $2m+2$, and the game's
  question distribution is undefined at that dimension for the admissible
  parameters: it requires $2m+2$ to divide $q$, while admissibility supplies
  only $m \mid q$; the repair is a statement-level reformulation of an
  imported theorem, not a new proof.  Second, the import of the classical
  game's quantum soundness from the tensor-code paper carries two open
  verification obligations at every dimension: the asserted identity with
  the cited tensor-code test is not an identity and must be replaced by a
  reduction with a value bound, and the source's instantiation of the
  theorem's auxiliary parameter violates its lower bound at small
  dimensions (Section~\ref{sec:import-obligations}).
  \item \textbf{Estimated weight:} for the recommended direction, about
  \textbf{3--5 new definitions or interfaces} (a line-point distribution with
  directly sampled index at arbitrary dimension, the directly indexed game,
  and a restated soundness import) plus \textbf{7--10 supporting lemmas},
  including the strategy-transport layer between the directly indexed game
  and the $\chi$-based game (parsing pullback, seed-fiber averaging with
  its open synchronicity-branch control, and a dilation--compression round
  trip of the soundness conclusion); several hundred lines of formal proof.
  Discharging the import obligations of
  Section~\ref{sec:import-obligations} adds the pair-measurement
  construction with its distributional matching, the branch-weight value
  bound, and the parameter-bound case split.  The padding fallback is
  heavier, roughly 6--10 reworked lemmas.
  \item \textbf{Mathlib/project split:} approximately \textbf{15\% Mathlib}
  and \textbf{85\% project-local}.  Mathlib supplies uniform distributions on
  finite types and pushforward identities; the substance is choosing the
  statement of the imported soundness theorem and rechecking the combining
  subtree against it.
  \item \textbf{Key Mathlib inputs:} uniform distributions on finite types
  and their maps, and elementary divisibility arithmetic on powers of two.
  The main inputs are project-local: the classical game, its conditionally
  linear question distribution, and the imported soundness statement.
  Tracked in \localissue{0002}.
\end{itemize}

\section{Key theorem forms}

Throughout, $q = 2^k$ with $k$ odd is an admissible field size, the tuple
$(q,m,d)$ is admissible when moreover $d \ge 1$ and $m \mid q$, and
$\mathrm{ideg}_{d,m}(\F_q)$ denotes the $m$-variate polynomial representatives
over $\F_q$ of individual degree at most $d$.  The comparison involves four
theorem forms.
\begin{enumerate}
  \item \textbf{Source theorem form} (the polynomial-pair measurement).  There
  is a function $\delta_S(\eps,m,d,q) = a(md)^a(\eps^b + q^{-b} + 2^{-bmd})$,
  for universal constants $a > 1$ and $0 < b < 1$, such that for every
  admissible $(q,m,d)$ and every projective strategy for the Pauli basis test
  succeeding with probability at least $1-\eps$, there is a projective
  measurement $\{\hat S_{g_X,g_Z}\}$ with outcomes pairs of polynomials in
  $\mathrm{ideg}_{d,m}(\F_q)$ whose evaluations are $\delta_S$-consistent, on
  average over a uniform point $u \in \F_q^m$ and in both bases
  $W \in \{X,Z\}$, with the corresponding points measurements
  $\{\hat M^{(\mathrm{Point},W),u}_a\}$ of the strategy, together with the
  symmetric counterpart of this relation.
  \item \textbf{Missing instantiation obligation.}  The classical low
  individual degree game with parameter tuple $(q,\, 2m+2,\, d,\, 1)$ is well
  defined, and the combined point and line measurements assembled in the
  analysis constitute a strategy for it of value at least
  $1 - \delta_{\mathrm{combine}}$.  The game's question distribution draws
  its axis and diagonal indices through an index function
  $\chi_{2m+2}\colon \F_q \to \{1,\ldots,2m+2\}$ with classes of equal size
  $q/(2m+2)$, so the game exists only when $2m+2$ divides $q$.
  \item \textbf{Restricted form supported by the present route.}  The source
  theorem form under the additional hypothesis $2m+2 \mid q$.  Among
  admissible tuples this hypothesis holds exactly when $m = 1$ and
  $q \ge 8$; it fails for every admissible tuple with $m \ge 2$ and for
  $(q,m) = (2,1)$.
  \item \textbf{Repaired import form} (recommended target).  Quantum
  soundness of the directly indexed variant of the classical low individual
  degree game over $\F_q^{m'}$, for every dimension $m' \ge 1$ and
  admissible $q$, with no divisibility hypothesis: the question
  distribution samples the axis index, and the diagonal prefix index,
  uniformly from $\{1,\ldots,m'\}$ directly rather than through an index
  function on $\F_q$, and line questions carry the sampled index itself in
  place of the seed $s \in \F_q$; the error is
  $\delta_{\mathrm{ld}}(\eps,q,m',d,k) = a(dm'k)^a(\eps^b + q^{-b} +
  2^{-bm'd})$.  This is a candidate target form, not an established one:
  proving it carries the same two import obligations as the $\chi$-based
  import (Section~\ref{sec:import-obligations}), and deriving the present
  form from it when $m' \mid q$ requires the strategy transport described
  under the candidate repairs below, whose synchronicity step is an open
  obligation.
\end{enumerate}

\section{The source assertion and the obstruction}

We examine the final step of the subsection ``Applying the classical
low-degree test'' in the analysis of the Pauli basis test in the compression
paper of Ji, Natarajan, Vidick, Wright, and Yuen (MIP$^{*}$=RE,
arXiv:2001.04383), whose quantum soundness component builds on the low
individual degree test of \cite{Ji2020LowIndividualDegree}.%
\footnote{Tracked in \localissue{0002}.  In the mirror of the source
consulted here, the polynomial-pair lemma and its proof are at
\path{references/qpbt-paper/14_analysis_of_the_pauli_basis_test.tex},
lines~1267--1288; the sub-line lemma is at lines~1063--1116 and the combined
line lemma at lines~1020--1034 of the same file.  The game parametrization
with its divisibility, the index function, the soundness statement with its
proof, and admissibility are at
\path{references/qpbt-paper/08_classical_and_quantum_low_degree_tests.tex},
lines~33--35, 214--219, 413--458, and 958--961.}
There, the combined point measurements
$\{\hat Q^{x,z,\alpha,\beta}_c\}$ and combined line measurements
$\{\hat Q^{\ell}_f\}$, indexed by points of $\F_q^{2m+2}$ and lines in
$\F_q^{2m+2}$, are asserted to form a strategy for the classical low
individual degree game with parameters
\[
  \mathsf{ldparams} = (q,\, 2m+2,\, d,\, 1)
  \tag{\path{lem:qld-4-7}}
\]
of value $1 - \delta_{\mathrm{combine}}$, and the quantum soundness theorem
for that game is applied to produce the global polynomial-pair measurement.
The two lemmas feeding this step are stated for the line-point distribution
over $\F_q^{2m+2}$.

The game, however, is parametrized by tuples $(q,m',d,k)$ with $q$ an
admissible field size \emph{and $m'$ dividing $q$}.  The divisibility is not
incidental: the question distribution is a conditionally linear distribution
in which a uniform seed coordinate $s \in \F_q$, read as an integer in
$\{0,\ldots,q-1\}$, determines the axis or diagonal index through the unique
$\chi(s)$ with
\[
  s = (\chi(s)-1)\,\frac{q}{m'} + r, \qquad 0 \le r < \frac{q}{m'},
  \tag{\path{eq:chi-func}}
\]
so that the $m'$ classes $\chi^{-1}(i)$ partition $\F_q$ into sets of equal
size $q/m'$ and the index is uniform.  For $m' \nmid q$ neither the index
function nor the line-point distribution built from it is defined.

Admissibility of $(q,m,d)$ supplies only $m \mid q$.  Since an admissible $q$
is a power of two, $m \mid q$ forces $m = 2^j$.  For $j \ge 1$ the integer
$m+1 \ge 3$ is odd, so $2m+2 = 2(m+1)$ has an odd factor exceeding $1$ and
divides no power of two; for $m = 1$ the condition $4 \mid q$ holds for every
admissible $q$ except $q = 2$.  Thus the instantiation at dimension $2m+2$ is
undefined for every admissible tuple with $m \ge 2$: for instance
$(q,m,d) = (8,2,1)$ is admissible, while $6 \nmid 8$.  The source performs
the instantiation without addressing the divisibility, so its proof route
establishes the polynomial-pair lemma, and with it the soundness theorem for
the Pauli basis test that rests on it, only for $m = 1$ and $q \ge 8$.

\section{What the cited soundness theorem requires: the import obligations}
\label{sec:import-obligations}

The quantum soundness theorem for the classical game inherits the constraint
$m' \mid q$ solely because it is stated for the game just described; its
proof fixes an admissible $q$ with $m' \mid q$, asserts that the game with
simultaneity parameter $1$ is \emph{identical} to the two-prover tensor-code
test for the Reed--Solomon tensor code $\mathcal C^{\ot m'}$ of the same
authors' \emph{Quantum soundness of testing tensor codes} (arXiv:2111.08131)
\cite{Ji2021Quantum}, and concludes by Theorem~4.7 there, instantiated with
that theorem's auxiliary integer parameter --- denoted $k$ in that source,
and written $K$ here to avoid clashing with the field exponent and the
simultaneity parameter --- chosen as $(m')^3 d$.  Neither
the asserted identity nor the instantiation is complete as stated, so the
import carries two verification obligations, which this section details and
which the blueprint's proof node for the soundness statement records as
open.  Both obligations are independent of the divisibility question: they
arise at every dimension, for the $\chi$-based game where it is defined and
equally for the directly indexed variant proposed under the candidate
repairs.

\paragraph{The cited test.}
In the two-prover tensor-code test (Figures~1 and~2 of
\cite{Ji2021Quantum}; the two-prover strategy format and the goodness
measure are Definitions~4.5 and~4.6 there), the referee performs one of
three subtests with probability $1/3$ each, assigning the two roles of each
subtest to the two provers uniformly at random.  \emph{Axis-parallel lines
test}: sample a uniform point $u \in [n]^{m'}$ and a uniform axis
$j \in \{1,\ldots,m'\}$; send the axis-parallel line $\ell$ through $u$ in
direction $j$ to one prover, who answers a codeword $g \in \mathcal C$;
send $u$ to the other, who answers a field value $a$; accept iff
$g(u_j) = a$.  \emph{Subcube commutation test}: sample
$j \in \{1,\ldots,m'\}$ and a subcube $H \subseteq [n]^{m'}$ by fixing the
last $j-1$ coordinates uniformly; sample $u, v$ independently and uniformly
in $H$ and a uniform order bit; send the ordered pair $(u,v)$ or $(v,u)$ to
one prover, who answers a pair $(b_u, b_v)$ of field values; send $u$ to
the other, who answers $a$; accept iff $b_u = a$.  \emph{Synchronicity
test}: send one identical question --- a line, a point, or an ordered
pair, with probability $1/3$ each --- to both provers and accept iff their
answers agree.  Theorem~4.7 supposes a two-prover strategy passing this
test with probability $1-\eps$ together with an integer parameter
$K \ge 12 m' t$, where $t$ is the interpolation parameter of the base code
--- for the degree-$d$ Reed--Solomon code, of blocklength $n = q$ and
distance $q-d$, one has $t = d+1$ --- and produces projective measurements
$\{G_c\}$ and $\{\tilde G_c\}$ with outcomes in $\mathcal C^{\ot m'}$,
consistent with one prover's points measurements, with the other prover's
lines measurements, and with each other, up to
$\eta'(m',t,K,\eps,n^{-1}) = \mathrm{poly}(m',t,K) \cdot
\mathrm{poly}(\eps, n^{-1}, e^{-K/(m')^2})$.

\paragraph{Obligation 1: the game correspondence.}
The source asserts that the $\chi$-based game at parameters $(q,m',d,1)$
is this test.  It is not, in four respects.
(i)~The game has no pair questions.  A pair question of the test reveals
two points of a \emph{hidden} subcube --- the subcube itself is not sent
--- and is answered by two field values, while the game's diagonal-line
question reveals a full line description and is answered by a degree-$m'd$
polynomial.  A reduction must construct the pair measurements from the
diagonal-line measurements by post-processing (measure a line through $u$
and $v$ and evaluate the answered polynomial at the two parameters), and
both directions of the information mismatch matter: the game's line
description carries an index and a seed, neither of which is a function of
the pair $(u,v)$, so the constructed measurement must average over the
compatible indices and seeds; and the game's diagonal lines fix the
\emph{first} $i-1$ coordinates where the test's subcubes fix the
\emph{last} $j-1$, so the coordinate conventions must be aligned.  The
construction must preserve projectivity and must match the law of the
hidden-subcube pair with the law induced by the game's diagonal sampling.
(ii)~The game's line questions carry the seed $s \in \F_q$; the test's
carry only the line.  A strategy may depend on the seed beyond its class
$\chi^{-1}(i)$, and removing this dependence runs into the fiber-averaging
difficulty made explicit under the candidate repairs: averaging preserves
the line-versus-point branches exactly but not the synchronicity branches.
(iii)~The branch weights differ: each nontrivial cross-test has mass $1/3$
in the test against $2/9$ in the game.  Once the per-branch predicates are
matched, the intended bookkeeping bound is that the constructed strategy's
rejection probability in the test is at most $\tfrac32$ times the game's
rejection probability, up to the $O(1/(m'q))$ of item~(iv); this bound is
the target of the obligation, not an established fact.
(iv)~The game's predicate fixes an acceptance convention on diagonal
questions with zero direction --- an event of probability at most
$2/(m'q)$, recorded in the blueprint's remark on zero directions --- for
which the test has no counterpart branch.
Finally, the displayed conclusions of Theorem~4.7 pair one prover's points
measurements with the other prover's tensor-codeword measurement, and the
first prover's tensor-codeword measurement with the second prover's lines
measurements, while the game-level statement asserts the point consistency
in both directions; deriving the stated form uses the role symmetrization
of the test and belongs to the correspondence.  None of these steps is
difficult to state, but none is carried out in the source, and together
they replace the asserted one-line identity by a reduction carrying its
own value bound.

\paragraph{Obligation 2: the parameter bound.}
Theorem~4.7 requires $K \ge 12 m' t = 12 m' (d+1)$.  The source
instantiates $K = (m')^3 d$, which meets the bound iff
$(m')^2 d \ge 12(d+1)$: this holds for all $d \ge 1$ when $m' \ge 5$,
holds at $m' = 4$ exactly when $d \ge 3$, and fails for every $d$ when
$m' \le 3$.  In the game's parameter domain the dimension divides the
admissible field size and is therefore a power of two, so the source's
instantiation violates the hypothesis exactly at $m' \in \{1,2\}$ for
every $d$ and at $m' = 4$ for $d \in \{1,2\}$ --- in particular, in the
multilinear regime $d = 1$, at every dimension $m' \le 4$.  A candidate
repair is to instantiate $K = \max((m')^3 d,\, 12 m'(d+1))$; the open
obligation is then to re-absorb the error, that is, to check that
$\mathrm{poly}(m',t,K)$ and $e^{-K/(m')^2}$ still collapse to the form
$a(dm')^a(\eps^b + q^{-b} + 2^{-bm'd})$ in the small-$m'$ cases where the
maximum is $12m'(d+1)$ --- plausible, since there
$e^{-K/(m')^2} = e^{-12(d+1)/m'}$ decays in $d$ at a rate comparable to
$2^{-m'd}$ for $m' \le 4$, but not checked here.

In the correspondence of Obligation~1, the only role of the divisibility
$m' \mid q$ is to make the $\chi$-sampled index uniform.  The underlying
mathematics does not need the divisibility.  The low
individual degree test of \cite{Ji2020LowIndividualDegree} samples its axis
index $i$ uniformly from $\{1,\ldots,m'\}$ directly, with no index function
on $\F_q$ and no divisibility hypothesis,%
\footnote{See the test figure in \path{references/ldt-paper/test_definition.tex}.}
and the tensor-code test is likewise a test about the code
$\mathcal C^{\ot m'}$, not about an encoding of indices into $\F_q$.  The
$\chi$-based packaging exists so that the question distribution is
conditionally linear, which the introspection machinery requires for the
game the players actually play --- the Pauli basis test at dimension $m$,
where admissibility does supply the divisibility.  The game at dimension
$2m+2$ never appears in a verifier; it is an internal analysis device in the
proof of the polynomial-pair lemma, so nothing there requires its question
distribution to be conditionally linear.

\section{Candidate repairs}

\paragraph{Direct index sampling (candidate direction, recommended).}
Define, for arbitrary $m' \ge 1$, the line-point distribution over
$\F_q^{m'}$ that samples the uniform point, the uniform axis index
$i \in \{1,\ldots,m'\}$, and the diagonal direction with vanishing first
$i-1$ coordinates, directly; define the \emph{directly indexed game} with
this question distribution, whose line questions carry the index $i$ (and,
for diagonal lines, the projected direction) in place of the seed
$s \in \F_q$; and state the imported soundness theorem for that game
(form~4 above).  This removes the divisibility obstruction and nothing
else: the directly indexed game still asks for polynomials on revealed
lines, has no pair questions, and weighs its branches as the present game
does, so establishing form~4 carries the same two import obligations as
the $\chi$-based import (Section~\ref{sec:import-obligations}), now at
arbitrary dimension.

The relation to the present game when $m' \mid q$ must moreover be proved
at the level of games, not merely of question distributions, because the
two games have different question alphabets: an $\mathsf{ALine}$ question
of the $\chi$-based game is a pair $(u_0, s)$ and a $\mathsf{DLine}$
question a triple $(u_0, s, v')$, so a strategy may select different
measurements for seeds $s \ne s'$ with $\chi(s) = \chi(s')$, and equality
of the induced laws of the geometric data does not by itself transport
such a strategy.  The transport layer decomposes into the following
obligations, none of which is carried out here.
\begin{itemize}
  \item \emph{Distribution lemma.}  When $m' \mid q$, the image of the
  conditionally linear question distribution under the parsing map
  $(u_0,s) \mapsto (u_0,\chi(s))$, $(u_0,s,v') \mapsto (u_0,\chi(s),v')$
  is the direct distribution, and conditioned on the parsed question the
  seed is uniform on the class $\chi^{-1}(i)$, independently of everything
  else.
  \item \emph{Parsing pullback.}  A strategy for the directly indexed game
  composed with the parsing map is a seed-oblivious strategy for the
  $\chi$-based game.  Since the win predicate depends on the question only
  through the parsed data, and a seed-oblivious strategy answers
  identically across each fiber --- in particular on the synchronicity
  branches, where both players hold the same seed --- the value should be
  preserved exactly on every branch, and projectivity is preserved.  This
  is the direction the combining subtree needs.
  \item \emph{Seed-fiber averaging: what it preserves and what it does
  not.}  In the opposite direction, average each line measurement over its
  seed fiber: attach to a direct question $(u_0,i)$, respectively
  $(u_0,i,v')$, the average over uniform $s \in \chi^{-1}(i)$ of the
  measurements attached to $(u_0,s)$, respectively $(u_0,s,v')$, leaving
  the points measurements unchanged.  By the conditional uniformity of the
  seed this preserves the value of every \emph{line-versus-point} branch
  exactly: there the paired question carries no seed, so the branch value
  is an average over the fiber of per-seed values, which the averaged
  measurement computes.  It does \emph{not} preserve the synchronicity
  branches.  There both players hold the same seed, so the $\chi$-based
  branch value is the diagonal average
  $\E_{s}\sum_f \langle M^{(u_0,s)}_f \ot \widetilde M^{(u_0,s)}_f
  \rangle$, while the averaged strategy realizes the double average
  $\E_{s,s'}\sum_f \langle M^{(u_0,s)}_f \ot \widetilde M^{(u_0,s')}_f
  \rangle$ over independent seeds in the fiber; the two differ whenever
  the players' measurements vary across a fiber in a correlated way (if
  both apply a common seed-dependent relabelling, every diagonal term is
  large and the cross terms are not), and neither quantity dominates the
  other in general.  An earlier revision of this note claimed that the
  averaging preserves the exact value, and with it synchronicity; that
  claim was wrong and is withdrawn.  What needs proof is control of the
  synchronicity branches under de-seeding.  Candidate approaches: show
  that a strategy of value $1-\eps$ in the $\chi$-based game is, with
  polynomial loss, close to a seed-oblivious one; or state form~4 with a
  synchronicity hypothesis in the doubled-average form and re-derive it on
  the tensor-code side.  Choosing and proving one of these is an open
  obligation.
  \item \emph{Projectivity round trip.}  The averaged strategy is a POVM
  strategy even when the given one is projective, while the imported
  theorem, like the present one, takes projective strategies.  The
  intended round trip is: dilate the averaged strategy to a projective one
  (Naimark preserves all product correlations
  $\langle X \ot Y \rangle$ of the given strategy against the dilated
  state --- hence the value of every branch the averaging preserved ---
  but nothing restores the diagonal seed correlations once the averaging
  has been performed); apply form~4; and compress the resulting polynomial
  measurement along the dilation isometry back to a POVM in the required
  measurement class on the original space.  That this compression
  preserves the consistency conclusions must itself be proved; it is
  stated here as an obligation, not established.  (Alternatively, form~4
  may be stated for POVM strategies outright, moving the dilation inside
  the import.)
\end{itemize}

The combining subtree would then be reread with the direct distribution
and the directly indexed game at dimension $2m+2$: every statement becomes
well defined for all admissible tuples, and the expectation is that the
proofs survive unchanged, because the measurements constructed there are
indexed by the geometric data --- points, lines, and polynomials on them
--- and so constitute a strategy for the directly indexed game outright,
with no fiber averaging.  Verifying, lemma by lemma, that those proofs use
only the laws of the sampled objects and never the conditional linearity
of the sampling is part of this direction's obligations.  This follows the
subtree protocol for a definition-level mismatch: the definition is
corrected at the lowest point rather than each downstream lemma acquiring
a hypothesis.  The transport layer costs one parsing-map definition, the
distribution lemma, the pullback lemma, whichever synchronicity control is
chosen, and the dilation--compression lemma; the first three are an
estimated one to two hundred lines on top of the restated import, with the
blueprint's existing Naimark interface supplying the dilation, while the
synchronicity control is the step whose weight is genuinely uncertain.
The external check that the tensor-code test of arXiv:2111.08131 samples
its axes and subcubes directly has been made against that source: Figure~1
there draws the axis index and the subcube data uniformly and directly,
with no index encoding into $\F_q$.

\paragraph{Padding the dimension.}
Alternatively, embed the combined polynomial into
$\mathrm{ideg}_{d,\bar m}(\F_q)$ for $\bar m = 4m$, the least power of two
with $\bar m \ge 2m+2$, ignoring the $2m-2$ dummy variables.  Then
$\bar m \mid q$ holds whenever $q \ge 4m$, and the existing game applies at
dimension $\bar m$.  In the complementary regime $m \in \{q/2,\, q\}$ the
conclusion is vacuous: $md \ge q/2$ gives
$\delta_S \ge a(md)^a q^{-b} \ge a\,2^{-a} q^{a-b} \ge a\,2^{-b} \ge 1$
for $a \ge 2$, while the consistency quantities bounded by the theorem never
exceed $1$.  This route stays inside the present definitions but reworks the
sub-line distribution (a fourth coordinate block of dummies), the combining
map, and the wrong-block Schwartz--Zippel analysis, and it moves the
combining lemmas away from their source statements at dimension $2m+2$.

\paragraph{Unequal index classes.}
One could instead define an index function at dimension $2m+2$ with classes
of sizes $\lfloor q/(2m+2)\rfloor$ and $\lceil q/(2m+2)\rceil$.  The index
is then only approximately uniform, the identification with the external
test fails exactly, and every distributional identity in the sub-line
analysis acquires an error term; this amounts to re-proving the imported
soundness theorem for a perturbed game and is not recommended.

\section{Comparison with the blueprint}

The blueprint \cite{MIPStarREBlueprint} keeps the sub-line lemma, the
combined line lemma, and the polynomial-pair lemma in their source-shaped
statements, with no divisibility hypothesis.  Their proofs retain the source
constructions and mark the unresolved obligation at the exact steps: the
first sampling step of the sub-line construction, where $\chi_{2m+2}$ and
the line-point distribution over $\F_q^{2m+2}$ are invoked, and the
instantiation step of the polynomial-pair proof, where $(q,2m+2,d,1)$ is
used as a game parameter tuple.  A chapter remark records the divisibility
arithmetic and cites this note; the final soundness theorem for the Pauli
basis test remains open modulo this obligation.  The blueprint's proof node
for the classical soundness import itself states the source's route --- the
identification with the two-prover tensor-code test and the application of
Theorem~4.7 --- and enumerates the two obligations of
Section~\ref{sec:import-obligations} as open, citing this note for the
detailed comparison and the case arithmetic.  An earlier revision had
instead added the hypothesis $2m+2 \mid q$ to the three lemmas; that was
statement drift in the sense of the proof-gap protocol, since the hypothesis
excludes all admissible $m \ge 2$ and does not appear in the source.

\section{Conclusion}

There is a theorem-level gap in the source proof route: the classical low
individual degree game invoked at dimension $2m+2$ is undefined for every
admissible tuple with $m \ge 2$ (and for $(q,m) = (2,1)$), because its
question distribution requires $2m+2 \mid q$ while admissibility supplies
only $m \mid q$.  The cited route therefore does not establish the
polynomial-pair lemma, nor the soundness of the Pauli basis test that
depends on it, beyond $m = 1$ and $q \ge 8$.  Independently of the
divisibility, the import of the classical game's quantum soundness carries
the two open verification obligations of
Section~\ref{sec:import-obligations}: the asserted identity with the
two-prover tensor-code test must be replaced by a reduction carrying a
value bound, and the source's instantiation of the auxiliary parameter
violates the theorem's lower bound at the small dimensions listed there.
We do not assert that any conclusion is false: the underlying tensor-code
soundness is insensitive to how the axis index is sampled, and both import
obligations concern verification, not counterexample.  For the divisibility
repair, the recommended candidate direction is direct index sampling ---
restate the imported soundness theorem for the directly indexed game at
arbitrary dimension; prove the distribution and parsing-pullback lemmas;
settle the synchronicity control for de-seeding together with the
dilation--compression round trip, so that the present soundness statement
would follow from the repaired import; and reread the combining subtree
against the directly indexed game --- a direction whose open obligations
are enumerated under the candidate repairs, not an established route, with
dimension padding as the fallback confined to the present game definition.

\bibliographystyle{alpha}
\bibliography{references}

\end{document}
